\documentclass[11pt,a4paper]{article}
\usepackage[margin=24mm]{geometry}
\usepackage{amsmath,amssymb,booktabs,enumitem}
\usepackage[T1]{fontenc}
\usepackage{lmodern}
\usepackage[hidelinks]{hyperref}
\setlength{\parindent}{0pt}
\setlength{\parskip}{5pt}
\setlist{nosep,leftmargin=*}
\newcommand{\C}{\mathbb C}
\newcommand{\E}{\mathcal E}
\newcommand{\add}{\operatorname{add}}
\newcommand{\Ext}{\operatorname{Ext}}
\title{Problem 00000000567: the ordinary representation category\texorpdfstring{\\}{ }already fails in type \texorpdfstring{$A_1$}{A1}}
\author{\texttt{lizaixi01}}
\date{24 October 2026}
\begin{document}
\begin{center}
{\Large\bfseries Problem 00000000567: the ordinary representation category\\
already fails in type $A_1$}\par
\smallskip
\texttt{lizaixi01}\quad -- \quad 24 October 2026
\end{center}

\section*{Statement and scope}
The first clause asserts that the number of pairwise nonisomorphic maximal
Ext-orthogonal objects in the idempotent completion of the representation
category of a Dynkin quiver is
\[
  \prod_{i=1}^{n}\frac{e_i+h+1}{e_i+1}.
\]
Take the ordinary category of finite-dimensional complex representations of the
one-vertex, no-arrow quiver $A_1$. Let $\E$ be its Karoubi envelope. We prove
two incompatible counts: with all object isomorphism classes there are
infinitely many; with the customary additional \emph{basic} convention there
is exactly one. The displayed product is $2$.

We use the standard additive meaning of maximal rigidity, rather than a strict
ordering of objects that would forbid repeated summands. Write $X\in\add T$
when $X$ is a direct summand of a finite direct sum of copies of $T$.
An object $T$ is maximal rigid if
\[
 \Ext^1(T,T)=0,\qquad
 \bigl(\Ext^1(U,U)=0,\ T\in\add U\bigr)\Longrightarrow U\in\add T.
\]
We also prove the stronger, two-sided cluster-tilting condition
\[
 \add T=\{X:\Ext^1(T,X)=0\}=\{X:\Ext^1(X,T)=0\}.
\]
These are the add-closure conditions in Burban--Iyama--Keller--Reiten,
Definition~1.1~\cite{bikr}. Basic objects are finite direct sums of pairwise
nonisomorphic nonzero indecomposables.

This disproof concerns the printed \emph{ordinary representation category}.
It makes no assertion against the usual Catalan enumeration in the derived
orbit cluster category $D^b(\operatorname{mod}H)/(\tau^{-1}[1])$, which is a
different category~\cite{bmrrt}. The preprojective-algebra clause is not used.

\section*{The complete category and its idempotents}
A representation of this quiver consists of one finite-dimensional vector
space: there are no arrow maps. Evaluation at the vertex and the constant
representation functor are inverse equivalences. In Lean the category is
literally the functor category from the native path category of the actual
arrowless quiver to \texttt{FGModuleCat} over $\C$.

For an arbitrary Karoubi object $(V,p)$ with $p^2=p$, evaluation yields an actual
linear idempotent on $V$. Its range is finite dimensional. The maps
\[
 V\xrightarrow{\,p\,}\operatorname{im}p
   \xrightarrow{\,\iota\,}V
\]
satisfy $\iota p=p$ on $V$ and $p\iota=1$ on $\operatorname{im}p$ (composition
is written in function order here). Their representation-level lifts give an
explicit isomorphism of $(V,p)$ with the identity idempotent on
$\operatorname{im}p$. Thus all objects of $\E$, including arbitrary formal
direct factors, are covered by the genuine equivalence
$\E\simeq\operatorname{Vect}_{\C,\mathrm{fd}}$.

\newpage
\section*{Actual Ext, finite sums, and classification}
Every epimorphism between finite-dimensional vector spaces is surjective and
admits the required linear lifts. Hence every object is projective.
Equivalence transports projectivity to every object of $\E$, and the native
abelian-category groups satisfy
\[
                    \Ext^1_{\E}(X,Y)=0\quad\hbox{for all }X,Y.
\]
These are actual derived-category Ext groups, not values assigned to a
replacement relation.

Let $d(X)$ be the actual finite dimension of the corresponding vector space.
A basis gives $X\simeq\C^{d(X)}$. Native finite biproducts satisfy
\[
 d\!\left(\bigoplus_i X_i\right)=\sum_i d(X_i).
\]
If $a\le b$, extension of coordinates by zero and restriction to the first $a$
coordinates give a genuine retraction $\C^a\longrightarrow\C^b\longrightarrow
\C^a$. Conversely any retraction induces an injective linear map, so it cannot
decrease the available target dimension.

If $T\ne0$ and $m=d(X)$, then $m\le m\,d(T)$. Applying the preceding
retraction to a finite biproduct of $m$ copies of $T$ gives $X\in\add T$.
If $T=0$, no such finite biproduct has the nonzero object $\C$ as a direct
summand. Consequently
\[
 T\text{ is maximal rigid}\quad\Longleftrightarrow\quad T\ne0
 \quad\Longleftrightarrow\quad T\text{ satisfies both Ext-orthogonal equalities}.
\]
In particular every actual $\C^{m+1}$ is a cluster-tilting object under these
conditions. Different dimensions prohibit isomorphisms. Taking the quotient of
the \emph{full} class of qualifying Karoubi objects by actual categorical
isomorphism therefore admits an injection from $\mathbb N$ and is infinite.
No finite-cardinality convention on an infinite type is used.

An object is indecomposable precisely when its dimension is one. Indeed a
positive dimension larger than one splits as $1+(d-1)$; a dimension-one
biproduct decomposition forces one summand to have dimension zero.
Therefore all nonzero indecomposables are isomorphic to $\C$. A finite
pairwise-nonisomorphic indecomposable decomposition has at most one summand.
Conversely the zero object and every dimension-one object have such a
decomposition. Thus
\[
 T\text{ basic}\Longleftrightarrow d(T)\le1,
 \qquad T\text{ basic and maximal rigid}\Longleftrightarrow d(T)=1.
\]
The quotient of all basic maximal-rigid objects is nonempty and a subsingleton;
its actual finite cardinality is exactly $1$.

\section*{Coxeter data derived from the root system}
The formalization constructs the native rank-one root pairing on the complex
line, with roots $1,-1$, coroot functionals $x\mapsto2x,-2x$, and evaluation
pairing. Both root families span. The singleton positive simple root is an
actual native base, with simple pairing $2$; its index set is explicitly
identified with the quiver vertex set. This is the complexified standard
$A_1$ root system.

\newpage
The Coxeter element is the product of its enumerated simple reflections. Here
the genuine reflection formula gives $s(x)=x-2x=-x$. Its square is the identity
and it is not the identity, so the actual group-theoretic order is $h=2$.
The scalar $-1$ is a primitive $h$-th root of unity. The full complex line is
the $-1$ eigenspace, and there are no other eigenvalues. Filtering the integers
$0\le e<h$ by the actual eigenvalue condition $(-1)^e$ gives exactly $\{1\}$,
with cardinality equal to the actual rank. Thus $e_1=1$ is derived from the
reflection representation. The rank-one data also agree with
Fomin--Reading, Section~13, printed page~29~\cite{fr}.
The conjectured expression is therefore the genuine product
\[
   \prod_{e\in\{1\}}\frac{h+e+1}{e+1}
       =\frac{2+1+1}{1+1}=2.
\]
It agrees with neither the infinite full-object count nor the basic count $1$.
This refutes the first necessary clause and hence the printed conjunction.

\section*{Formalization and reproducibility}
The complete argument is in \texttt{lean/Main.lean}, namespace
\texttt{TLMC567}. The table locates its main result and principal bridges.

\begin{center}
\small
\begin{tabular}{ll}
\toprule
Bridge & Declaration\\
\midrule
Main result & \texttt{complete\_counterexample}\\
All representations and their envelope & \texttt{envelopeEquivalence}\\
Arbitrary actual idempotent image & \texttt{explicitImageIso}\\
Native Ext vanishing & \texttt{envelope\_ext\_vanishes}\\
Standard maximality characterization & \texttt{maximalRigid\_iff\_nonzero}\\
Two-sided Ext characterization & \texttt{clusterTilting\_iff\_nonzero}\\
Full-object quotient infinity & \texttt{allClass\_injective}\\
Complete basic-object quotient count & \texttt{basic\_class\_count}\\
Actual Coxeter order and exponents & \texttt{coxeter\_number\_eq\_two}\\
 & \texttt{coxeter\_exponents\_eq}\\
\bottomrule
\end{tabular}
\end{center}

The project pins Lean~4.33.0 and Mathlib commit\par
\texttt{db584cd6d46c92f209a44c0f1c829460d327499d}.\par
First run \texttt{lean --version}; then use \texttt{lean/verify.sh} for the
root build, warning-as-error namespace audit and official full fresh-kernel replay.
The accompanying receipts bind exact sources, commands, exit status and
original execution data. \texttt{Check.lean} enumerates every owned declaration
and invokes ordinary \texttt{\#print axioms} on each, including generated names.\par
The allowed foundational axioms are only \texttt{propext},
\texttt{Classical.choice}, and \texttt{Quot.sound}.

\begin{thebibliography}{9}
\small\raggedright
\bibitem{bikr} I. Burban, O. Iyama, B. Keller and I. Reiten,
\emph{Cluster tilting for one-dimensional hypersurface singularities},
author manuscript, Definition~1.1, printed p.~3.
\url{https://webusers.imj-prg.fr/~bernhard.keller/publ/ClusterSing.pdf}.
\bibitem{bmrrt} A. B. Buan, R. Marsh, M. Reineke, I. Reiten and G. Todorov,
\emph{Tilting theory and cluster combinatorics},
arXiv:math/0402054v1 (2004), printed pp.~1--3, 9--11.
\url{https://arxiv.org/pdf/math/0402054v1}.
\bibitem{fr} S. Fomin and N. Reading,
\emph{Generalized cluster complexes and Coxeter combinatorics},
arXiv:math/0505085v3 (2006), printed pp.~28--29.
\url{https://arxiv.org/pdf/math/0505085v3}.
\end{thebibliography}
\end{document}
